% ----------------------------------------------------------------------
%  3.2. Основы электронной теории кристаллов; 3.2.1. Свободные электроны
% ----------------------------------------------------------------------
\begin{orientir}
Почему металл проводит ток и тепло, а алмаз~— нет? Ответ дают
электроны, но описывать их нужно квантово. Мы пройдём два шага:
сначала (3.2.1) забудем о решётке и рассмотрим «газ» свободных
электронов в ящике~— это модель Зоммерфельда; затем (3.2.2) учтём
периодическое поле ионов и получим \emph{энергетические зоны}~— теорию
Блоха и Бриллюэна. Ключевые понятия раздела: энергия и сфера Ферми,
плотность состояний, функция Ферми--Дирака, зоны Бриллюэна, поверхность
Ферми.
\end{orientir}

Металлы занимают особое место среди твёрдых тел: высокая электро- и
теплопроводность, металлический блеск и непрозрачность, ковкость,
особые термоэлектрические и гальваномагнитные свойства, сильный
парамагнетизм и ферромагнетизм. Электронная теория твёрдого тела
создавалась прежде всего для того, чтобы объяснить эти свойства; уже
потом, отталкиваясь от металлов, удалось понять и неметаллы.

Задолго до квантовой механики Друде предложил \emph{классическую}
модель: металл~— ящик с газом свободных электронов (подробнее
п.~3.5.1). Модель имела яркие успехи (закон Ома, закон Видемана--Франца,
п.~3.7.3) и серьёзные провалы. Классическая теория не смогла объяснить
температурную зависимость теплоёмкости и парамагнитной восприимчивости
электронов проводимости. Оставалось загадкой и главное: почему
плотная среда «прозрачна» для электронов и почему электроны проводимости
ведут себя как газ невзаимодействующих частиц. Ответы дала только
квантовая теория.

\subsection{Квантовая теория свободных электронов}\label{ssec:free-el}

\paragraph{Модель}
Теория Зоммерфельда делает три упрощения.
\begin{shagi}
\item Валентные электроны \emph{делокализованы} и движутся свободно, как
частицы газа; структура металла не учитывается.
\item Положительный заряд ионов «размазан» по объёму равномерно и
нейтрализует заряд электронов. Внутри металла потенциальная энергия
электрона постоянна, а у поверхности резко возрастает и не выпускает
электрон наружу.
\item Взаимное отталкивание электронов не учитывается (\term{приближение
независимых электронов}).
\end{shagi}
Металл превращается в \term{потенциальный ящик} (рис.~\ref{fig:pot-box}).
Глубина ящика $U_0$ для большинства металлов порядка 10~эВ. Энергию
удобно отсчитывать от дна ящика; тогда внутри металла энергия электрона
чисто кинетическая: $E=mv^2/2=p^2/2m$.

\begin{figure}[ht]
\centering
\begin{tikzpicture}[x=1cm, y=0.36cm]
  % стенки и дно ящика
  \draw[line width=1pt] (-1.2,10) -- (0,10) -- (0,0) -- (6,0) -- (6,10) -- (7.6,10);
  % заполненные состояния до уровня Ферми
  \fill[fmBlue!22] (0,0) rectangle (6,5.5);
  \foreach \y in {0.6,1.2,...,5.4}{\draw[fmBlue!60, line width=0.4pt] (0.1,\y) -- (5.9,\y);}
  \foreach \y in {6.0,6.6,...,9.6}{\draw[fmGray!50, line width=0.4pt, dashed] (0.1,\y) -- (5.9,\y);}
  \draw[fmRed, line width=1pt] (-0.2,5.5) -- (6.2,5.5) node[right, font=\small]{$\EF$};
  \node[font=\small, anchor=west] at (7.65,10) {уровень вакуума};
  \draw[dimline] (-0.6,0) -- node[left, font=\small]{$U_0$} (-0.6,10);
  \draw[dimline] (7.2,5.5) -- node[right, font=\small]{$\varphi$} (7.2,10);
  \draw[dimline] (0,-1.2) -- node[below, font=\small]{$L$} (6,-1.2);
  \node[font=\small, fmBlue!70!black, align=center] at (3,2.7) {заняты\\(при $T=0$)};
  \node[font=\small, fmGray, align=center] at (3,7.8) {свободны};
\end{tikzpicture}
\caption{Металл как потенциальный ящик размером $L$ и глубиной $U_0$.
Уровни энергии электронов квантованы (горизонтальные линии~— условно;
на самом деле они расположены неизмеримо гуще). При $T=0$ электроны
заполняют уровни снизу вверх до энергии Ферми $\EF$; $\varphi=U_0-\EF$~—
работа выхода}
\label{fig:pot-box}
\end{figure}

\paragraph{Электронный газ~— не классический газ}
Самый замечательный вывод теории: электронный газ при плотностях,
характерных для металлов, ведёт себя совсем не так, как обычный газ
молекул, подчиняющийся статистике Максвелла--Больцмана. Электроны имеют
полуцелый спин и подчиняются \term{принципу Паули}, а значит, квантовой
\term{статистике Ферми--Дирака}. Посмотрим, к чему это приводит.

\paragraph{Шаг 1: энергия квантуется}
Электрон~— волна, и его состояние в ящике описывает уравнение Шрёдингера.
Для одномерного ящика длиной $L$ (потенциальная энергия внутри равна
нулю):
\begin{equation}\label{eq:schr-1d}
-\frac{\hbar^2}{2m}\frac{\d^2\psi}{\d x^2}=E\psi .
\end{equation}
Электрон не может покинуть ящик, поэтому $\psi(0)=\psi(L)=0$. Этому
условию удовлетворяют \emph{стоячие волны}, у которых на длине $L$
укладывается целое число полуволн:
\[
\psi_n(x)=A\sin\frac{\pi n x}{L},\qquad n=1,2,3,\ldots
\]
Подставив $\psi_n$ в \eqref{eq:schr-1d}, получим разрешённые значения
энергии:
\begin{keyeq}\label{eq:box-levels}
E_n=\frac{\hbar^2\pi^2n^2}{2mL^2}.
\end{keyeq}
Энергия может принимать только эти, \emph{собственные}, значения~—
спектр дискретен. Но для куска металла уровни расположены фантастически
густо: при $L=1$~см и $\hbar^2/2m=0{,}0381~\text{эВ}\cdot\text{нм}^2$
имеем $E_1=0{,}0381\pi^2/(10^7)^2\approx4\cdot10^{-15}$~эВ. Поэтому
спектр называют \term{квазинепрерывным}.

\paragraph{Шаг 2: состояния в пространстве волновых векторов}
Для кристалла удобнее не нулевые, а \term{периодические граничные
условия}: волновая функция повторяется через длину $L$ по каждой оси,
$\psi(x+L,y,z)=\psi(x,y,z)$ и~т.\,д. Тогда решения~— \emph{бегущие}
плоские волны
\[
\psi_{\vb k}(\vb r)=A\,e^{i\vb k\vb r},\qquad E=\frac{\hbar^2k^2}{2m},
\qquad \vb p=\hbar\vb k ,
\]
где $\vb k$~— \term{волновой вектор}. Условие периодичности требует
$e^{ik_xL}=1$, т.\,е.
\begin{equation}\label{eq:k-quant}
k_x=\frac{2\pi}{L}n_x,\quad k_y=\frac{2\pi}{L}n_y,\quad k_z=\frac{2\pi}{L}n_z,
\qquad n_x,n_y,n_z=0,\pm1,\pm2,\ldots
\end{equation}
Разрешённые $\vb k$ образуют в $\vb k$-пространстве кубическую сетку с
шагом $2\pi/L$ (рис.~\ref{fig:fermi-sphere}). На каждую точку сетки
приходится объём $(2\pi/L)^3=8\pi^3/V$, где $V=L^3$~— объём металла, и в
каждой точке по принципу Паули могут находиться \emph{два} электрона с
противоположными спинами.

\begin{smysl}[Импульсы или волновые векторы?]
В оригинале учебника рассуждение ведётся в пространстве импульсов: из
соотношения неопределённостей одно состояние занимает в фазовом
пространстве объём $h^3$, а в пространстве импульсов (при размере
металла $V$)~— объём $h^3/V$. Это то же самое: так как $\vb p=\hbar\vb
k$, объём $(2\pi/L)^3$ в $\vb k$-пространстве соответствует объёму
$(2\pi\hbar)^3/V=h^3/V$ в пространстве импульсов.
\end{smysl}

\paragraph{Шаг 3: сфера Ферми}
При $T=0$ электроны занимают состояния с наименьшей энергией. Энергия
$E=\hbar^2k^2/2m$ зависит только от длины $\vb k$, поэтому занятые
состояния заполняют в $\vb k$-пространстве \term{сферу Ферми} радиуса
$\kF$. Её объём, делённый на объём одного состояния и умноженный на 2
(спин), равен числу электронов $N$:
\[
2\cdot\frac{\frac43\pi\kF^3}{8\pi^3/V}=N
\quad\Rightarrow\quad
\kF=\left(3\pi^2n\right)^{1/3},\qquad n=\frac NV .
\]
Радиус сферы Ферми зависит только от концентрации электронов $n$, а не
от размеров куска металла. Наибольшая энергия электронов при $T=0$~—
\term{энергия Ферми}:
\begin{keyeq}\label{eq:EF}
\EF=\frac{\hbar^2\kF^2}{2m}=\frac{\hbar^2}{2m}\left(3\pi^2n\right)^{2/3}.
\end{keyeq}
В оригинале эта формула записана через импульс:
$\EF=\frac{h^2}{2m}\left(\frac{3N}{8\pi V}\right)^{2/3}$~— это то же
самое.

\begin{figure}[ht]
\centering
\begin{tikzpicture}[scale=0.42]
  \def\kf{4.25}
  \fill[fmBlue!10] (0,0) circle (\kf);
  \foreach \i in {-6,...,6}{\foreach \j in {-6,...,6}{
    \pgfmathsetmacro{\rr}{sqrt(\i*\i+\j*\j)}
    \ifdim\rr pt<\kf pt
      \fill[fmBlue] (\i,\j) circle (0.17);
    \else
      \draw[fmGray, line width=0.4pt] (\i,\j) circle (0.15);
    \fi}}
  \draw[fmRed, line width=1pt] (0,0) circle (\kf);
  \draw[arr, fmRed] (0,0) -- (40:\kf) node[pos=0.62, above left=-1pt, font=\small]{$\kF$};
  \draw[-{Stealth[length=2mm]}, fmGray] (-7,0) -- (7.2,0) node[right, font=\small]{$k_x$};
  \draw[-{Stealth[length=2mm]}, fmGray] (0,-7) -- (0,7.2) node[above, font=\small]{$k_y$};
  \draw[dimline] (5,-6.6) -- node[below, font=\scriptsize]{$2\pi/L$} (6,-6.6);
  \node[font=\small, fmRed, anchor=west] at (6.5,4.2) {сфера Ферми};
  \draw[fmRed, line width=0.4pt] (6.55,4.0) -- (3.4,2.6);
\end{tikzpicture}
\caption{Сечение $\vb k$-пространства. Разрешённые волновые векторы
образуют сетку с шагом $2\pi/L$; в каждой точке~— два электрона. При
$T=0$ заняты все точки внутри сферы Ферми (синие) и свободны все точки
вне её. Для реального куска металла сетка неизмеримо гуще}
\label{fig:fermi-sphere}
\end{figure}

\begin{primer}[Электронный газ натрия]
Натрий одновалентен; плотность $\rho=0{,}97$~г/см$^3$, $A=23$. Концентрация
электронов равна концентрации атомов:
$n=\rho N_A/A=0{,}97\cdot6{,}02\cdot10^{23}/23=2{,}54\cdot10^{22}$~см$^{-3}=25{,}4$~нм$^{-3}$.
\begin{shagi}
\item $\kF=(3\pi^2\cdot25{,}4)^{1/3}=9{,}09$~нм$^{-1}$; длина волны
электрона на уровне Ферми $\lambda_\text{F}=2\pi/\kF=0{,}69$~нм.
\item $\EF=0{,}0381\cdot9{,}09^2=3{,}15$~эВ.
\item \term{Температура Ферми} $T_\text{F}=\EF/k=3{,}15/8{,}62\cdot10^{-5}\approx3{,}7\cdot10^4$~К.
\item Скорость Ферми $v_\text{F}=\hbar\kF/m\approx1{,}05\cdot10^6$~м/с~— в
десять раз больше средней тепловой скорости классической частицы при
комнатной температуре ($\sqrt{3kT/m}\approx1{,}2\cdot10^5$~м/с).
\end{shagi}
Даже при абсолютном нуле электроны в металле движутся с огромными
скоростями~— это прямое следствие принципа Паули. У классического газа
при $T=0$ все частицы покоились бы.
\end{primer}

Вычисленные по \eqref{eq:EF} энергии Ферми приведены в
табл.~\ref{tab:EF}. Для одновалентных металлов $n$ равна концентрации
атомов $N_A/V_a$ ($V_a=A/\rho$~— молярный объём); для двух- и
трёхвалентных её нужно умножить на~2 или~3.

\begin{table}[ht]
\centering
\caption{Энергии Ферми металлов в модели свободных электронов, эВ}
\label{tab:EF}
\small
\begin{tabular}{@{}cccccccc@{}}
\toprule
\multicolumn{8}{@{}l}{\emph{Одновалентные}} \\
Li & Na & K & Rb & Cs & Cu & Ag & Au \\
4,74 & 3,16 & 2,06 & 1,79 & 1,53 & 7,10 & 5,52 & 5,56 \\
\midrule
\multicolumn{5}{@{}l}{\emph{Двухвалентные}} & \multicolumn{3}{l}{\emph{Трёхвалентные}} \\
Be & Mg & Ca & Zn & Cd & Al & Ga & In \\
14,14 & 7,13 & 4,68 & 9,39 & 7,46 & 11,63 & 10,35 & 8,60 \\
\bottomrule
\end{tabular}
\end{table}

\paragraph{Проверка: рентгеновские спектры}
Энергию Ферми можно измерить. Если выбить электрон с внутренней оболочки
атома, на освободившееся место переходят электроны проводимости, и
металл излучает мягкие рентгеновские лучи. Энергия кванта зависит от
того, с какого уровня пришёл электрон, поэтому ширина полосы излучения
$\Delta E$ равна ширине занятой части спектра, т.\,е. $\EF$. Опыт хорошо
согласуется с табл.~\ref{tab:EF}.

Можно рассуждать и обратно~— по ширине полосы найти число свободных
электронов на атом $n_c$. Из \eqref{eq:EF} концентрация электронов
$n=(2m\Delta E/\hbar^2)^{3/2}/3\pi^2$, или в удобных единицах
\begin{equation}\label{eq:nc}
n_c=n\,\Omega=4{,}54\cdot\left(\frac{\Delta E}{1~\text{эВ}}\right)^{3/2}\cdot\frac{\Omega}{1~\text{нм}^3},
\end{equation}
где $\Omega=A/(\rho N_A)$~— объём, приходящийся на один атом. Для натрия
($\Delta E=3{,}15$~эВ, $\Omega=0{,}0394$~нм$^3$) получаем $n_c=1{,}0$~—
один свободный электрон на атом, как и должно быть.

\paragraph{Плотность состояний}
Сколько состояний приходится на единичный интервал энергии? Число
состояний (с учётом двух проекций спина) с энергией меньше $E$ в
единице объёма~— это число точек внутри сферы радиуса
$k=\sqrt{2mE}/\hbar$:
\[
N(E)=\frac{1}{3\pi^2}\left(\frac{2mE}{\hbar^2}\right)^{3/2}.
\]
Её производная~— \term{плотность состояний}:
\begin{keyeq}\label{eq:dos-free}
g(E)=\frac{\d N}{\d E}=\frac{1}{2\pi^2}\left(\frac{2m}{\hbar^2}\right)^{3/2}\sqrt E
=\frac{3n}{2\EF}\sqrt{\frac{E}{\EF}}.
\end{keyeq}
Плотность состояний растёт с энергией по параболическому закону $\sqrt
E$ (рис.~\ref{fig:fermi-dirac},\,\textit{б})\footnote{В оригинале
плотность распределения \emph{электронов} обозначена $n(E)$, а
плотность \emph{состояний} $N(E)$ отнесена к одному направлению спина,
так что $N(E)=n(E)/2$. Мы используем $g(E)$ с учётом обеих проекций
спина; при $T=0$ ниже $\EF$ она совпадает с $n(E)$ оригинала.}.

\paragraph{Средняя энергия}
При $T=0$ средняя энергия электрона
\[
\overline E=\frac{\int_0^{\EF}E\,g(E)\,\d E}{\int_0^{\EF}g(E)\,\d E}
=\frac{\int_0^{\EF}E^{3/2}\d E}{\int_0^{\EF}E^{1/2}\d E}
=\frac{\frac25\EF^{5/2}}{\frac23\EF^{3/2}}=\frac35\EF .
\]
Для натрия это около 1,9~эВ~— в 50 раз больше средней тепловой энергии
$\frac32kT=0{,}039$~эВ, которую имел бы классический газ при комнатной
температуре.

\paragraph{Распределение Ферми--Дирака}
При $T>0$ вероятность того, что состояние с энергией $E$ занято,
задаёт \term{функция Ферми--Дирака}:
\begin{keyeq}\label{eq:fermi-dirac}
f(E)=\frac{1}{\exp\!\left(\dfrac{E-\mu}{kT}\right)+1},
\end{keyeq}
где $\mu$~— химический потенциал. Состояния с $f=1$ заполнены целиком
(двумя электронами), с $f=\frac12$~— в среднем наполовину, с $f=0$~—
пусты. При $T=0$ химический потенциал равен энергии Ферми, $\mu=\EF$, и
функция становится «ступенькой»: $f=1$ при $E<\EF$ и $f=0$ при $E>\EF$.
Такой газ называют \term{полностью вырожденным}. С ростом температуры
химический потенциал чуть смещается:
\begin{equation}\label{eq:mu-T}
\mu=\EF\left[1-\frac{\pi^2}{12}\left(\frac{kT}{\EF}\right)^2\right],
\end{equation}
но для металлов $kT/\EF\sim0{,}01$, и этим смещением практически всегда
можно пренебречь.

\begin{figure}[ht]
\centering
\begin{subfigure}[b]{0.48\textwidth}\centering
\begin{tikzpicture}
\begin{axis}[width=\textwidth, height=5cm, xmin=0, xmax=1.6, ymin=0, ymax=1.15,
  xlabel={$E/\EF$}, ylabel={$f(E)$}, xtick={0,0.5,1,1.5}, ytick={0,0.5,1},
  legend style={font=\scriptsize, at={(0.98,0.95)}, anchor=north east, draw=none, fill=none}]
  \addplot[fmGray, thick, dashed] coordinates {(0,1) (1,1) (1,0) (1.6,0)};
  \addlegendentry{$T=0$}
  \addplot[fmRed, thick, domain=0:1.6, samples=200] {1/(exp(min(9,max(-9,(x-1)/0.06)))+1)};
  \addlegendentry{$T>0$}
  \draw[dimline] (axis cs:0.88,0.3) -- (axis cs:1.12,0.3);
  \node[font=\scriptsize, anchor=west] at (axis cs:1.13,0.3) {$\sim\!4kT$};
\end{axis}
\end{tikzpicture}
\caption{}
\end{subfigure}\hfill
\begin{subfigure}[b]{0.48\textwidth}\centering
\begin{tikzpicture}
\begin{axis}[width=\textwidth, height=5cm, xmin=0, xmax=1.6, ymin=0, ymax=1.45,
  xlabel={$E/\EF$}, ylabel={$g(E)$, $n(E)$}, xtick={0,0.5,1,1.5}, ytick=\empty]
  \addplot[fmBlue, thick, domain=0:1.6, samples=100] {sqrt(x)};
  \addplot[fill=fmBlue!25, draw=none, domain=0:1.6, samples=200] {sqrt(x)/(exp(min(9,max(-9,(x-1)/0.06)))+1)} \closedcycle;
  \addplot[fmRed, thick, domain=0:1.6, samples=200] {sqrt(x)/(exp(min(9,max(-9,(x-1)/0.06)))+1)};
  \node[font=\scriptsize, fmBlue, anchor=south east] at (axis cs:1.55,1.25) {$g(E)\propto\sqrt E$};
  \node[font=\scriptsize, fmRed, align=center] at (axis cs:0.55,0.35) {заняты:\\$n(E)=g(E)f(E)$};
\end{axis}
\end{tikzpicture}
\caption{}
\end{subfigure}
\caption{\textit{а}~— функция Ферми--Дирака при $T=0$ (ступенька) и
при $T>0$: размывается лишь узкая полоса шириной в несколько $kT$
около $\EF$ (масштаб по оси абсцисс для наглядности сильно растянут:
в действительности $kT/\EF\sim0{,}01$); \textit{б}~— плотность
состояний свободных электронов и распределение электронов по энергиям
с тепловым размытием около $\EF$}
\label{fig:fermi-dirac}
\end{figure}

\begin{glavnoe}
При нагреве \emph{не все} электроны могут увеличить энергию. Тепловая
энергия $kT$ даже при 1000~К составляет около 0,1~эВ, а $\EF\sim5$~эВ.
Электрон из глубины распределения мог бы перейти только на уровень,
который уже занят,~— это запрещает принцип Паули. Поэтому возбуждаются
лишь электроны в узкой полосе шириной $\sim kT$ около уровня Ферми~—
примерно $kT/\EF\sim1$\,\% всех электронов. Отсюда малая электронная
теплоёмкость металлов (п.~3.7.2) и слабый парамагнетизм электронов
проводимости (п.~3.6)~— то, чего не смогла объяснить классическая теория.
\end{glavnoe}

Вырождение снимается лишь при $kT\approx\EF$, т.\,е. при температурах
порядка десятков тысяч кельвинов, когда металл давно превратился в пар.
Тепловое размытие видно и на опыте: высокоэнергетический край полосы
мягкого рентгеновского излучения металлов размыт, и размытие растёт с
температурой, как и предсказывает теория.

\begin{itogi}
\begin{itemize}
\item Модель свободных электронов: металл~— потенциальный ящик, электроны
не взаимодействуют ни с ионами, ни друг с другом; хорошо работает для
щелочных металлов.
\item Энергия квантуется, но уровни расположены так густо, что спектр
квазинепрерывен. В $\vb k$-пространстве на одно состояние приходится
объём $8\pi^3/V$, в нём~— два электрона.
\item При $T=0$ электроны заполняют сферу Ферми радиуса
$\kF=(3\pi^2n)^{1/3}$; энергия Ферми $\EF=\hbar^2\kF^2/2m$ составляет
2--15~эВ, средняя энергия $\frac35\EF$.
\item Плотность состояний $g(E)\propto\sqrt E$.
\item При $T>0$ распределение Ферми--Дирака размывается лишь в полосе
$\sim kT$ около $\EF$; возбуждается около 1\,\% электронов.
\end{itemize}
\end{itogi}

\begin{proverka}
\item Почему уровни энергии электрона в металле можно считать
непрерывными, хотя они дискретны?
\item Во сколько раз изменится энергия Ферми, если концентрация
электронов вырастет в 8 раз?
\item Оцените энергию Ферми меди ($\rho=8{,}9$~г/см$^3$, $A=63{,}5$, один
свободный электрон на атом).
\item Почему средняя энергия электронов при $T=0$ не равна нулю?
\item Какая доля электронов натрия может поглотить тепловую энергию при
комнатной температуре?
\end{proverka}
